\rtmtight
\begin{tblr}{width=\textwidth,colspec={|Q[c,m,2.1cm]|X[l,m]|},hlines,vlines,hspan=minimal,rowsep=1.5pt,colsep=3pt,row{1}={bg=headgray,font=\bfseries}}
\SetCell{c,m}\hfil\logo\hfil & \textbf{POLITEKNIK NEGERI MALANG}\par\textbf{JURUSAN TEKNOLOGI INFORMASI}\par\textbf{PROGRAM STUDI: D4 TEKNIK INFORMATIKA} \\
\SetCell[c=2]{l} \textbf{RENCANA TUGAS MAHASISWA (RTM) DAN RUBRIK PENILAIAN} & \\
Nama Mata Kuliah & Bahasa Indonesia \\
Kode Mata Kuliah & RTI254003 \quad \textbf{sks / Jam perminggu:} 2 SKS/4 jam \quad \textbf{Semester:} 4 \\
Dosen Pengampu & Tim Pengajar Program Studi D4 TI \\
\end{tblr}
\assessmentblock{Tugas Laporan Teknis dan Penulisan}{Case Method dan penulisan akademik}{SCPMK0301-03001, SCPMK0401-03002}{Mahasiswa menyelesaikan tugas menulis laporan teknis, makalah, dan karya tulis ilmiah sederhana sesuai Sub-CPMK dan konteks mata kuliah.}{Menganalisis kebutuhan penulisan, menyusun kerangka, mengembangkan isi, menyunting, dan mendokumentasikan.}{Laporan teknis, makalah, atau karya tulis, dan/atau bahan presentasi.}{Ketepatan struktur dan sistematika tulisan & 10 & Tulisan memiliki struktur lengkap (pendahuluan, isi, penutup), sistematika runtut, dan setiap bagian berkontribusi pada tujuan penulisan. & Struktur tulisan sebagian besar runtut, tetapi beberapa bagian kurang lengkap atau tidak proporsional. & Tulisan tidak terstruktur, bagian penting hilang, atau sistematika tidak dapat diikuti. \\ \\
Penggunaan bahasa Indonesia baku dan EYD & 11 & Bahasa Indonesia baku digunakan secara konsisten, ejaan dan tanda baca sesuai EYD V, diksi tepat dan kalimat efektif. & Sebagian besar bahasa baku, tetapi masih ada kesalahan ejaan, tanda baca, atau diksi yang berulang. & Banyak kesalahan bahasa, ejaan tidak sesuai EYD, atau diksi tidak tepat sehingga mengganggu keterbacaan. \\ \\
Kedalaman isi dan relevansi konten & 9 & Isi tulisan relevan dengan topik, argumentasi kuat, didukung data/referensi yang tepat, dan analisis menunjukkan pemahaman mendalam. & Isi cukup relevan dan ada argumentasi, tetapi dukungan data atau kedalaman analisis masih terbatas. & Isi dangkal, tidak relevan, atau tidak didukung data/referensi yang memadai. \\}

\assessmentblock{Kuis 1 Ejaan dan Kalimat Efektif}{Kuis}{SCPMK0301-03001}{Kuis tertulis untuk mengukur pemahaman ejaan, tanda baca, diksi, dan kalimat efektif dalam bahasa Indonesia baku.}{Menjawab soal pilihan ganda dan/atau uraian singkat.}{Jawaban tertulis kuis.}{Penguasaan EYD dan tanda baca & 3 & Ejaan dan tanda baca diterapkan dengan benar dan konsisten pada semua contoh soal. & Sebagian besar ejaan dan tanda baca benar, tetapi ada beberapa kesalahan yang berulang. & Banyak kesalahan ejaan dan tanda baca menunjukkan pemahaman yang belum memadai. \\ \\
Ketepatan diksi dan kalimat efektif & 2 & Diksi tepat, kalimat efektif, lugas, dan tidak ambigu pada semua jawaban. & Diksi cukup tepat, tetapi beberapa kalimat masih rancu atau tidak efektif. & Diksi tidak tepat atau kalimat tidak efektif pada sebagian besar jawaban. \\}

\assessmentblock{UTS Laporan Teknis dan Tata Bahasa}{UTS berbasis penulisan}{SCPMK0301-03001, SCPMK0401-03002}{Evaluasi tengah semester mengukur kemampuan menulis laporan teknis dan menerapkan tata bahasa Indonesia yang benar.}{Menulis laporan teknis pendek atau merevisi teks sesuai soal UTS.}{Laporan teknis tertulis atau teks yang telah disunting.}{Kualitas laporan teknis & 8 & Laporan teknis terstruktur, isi relevan, bahasa baku, dan memenuhi semua komponen yang diminta. & Laporan cukup terstruktur, tetapi beberapa komponen kurang lengkap atau bahasa belum sepenuhnya baku. & Laporan tidak terstruktur, banyak komponen hilang, atau bahasa tidak baku secara signifikan. \\ \\
Penerapan tata bahasa dan sitasi & 7 & Tata bahasa benar, ejaan sesuai EYD V, dan sitasi (jika diminta) mengikuti format yang ditentukan. & Tata bahasa sebagian besar benar, tetapi ada kesalahan ejaan atau format sitasi yang tidak konsisten. & Banyak kesalahan tata bahasa atau sitasi tidak sesuai format yang ditentukan. \\}

\assessmentblock{Kuis 2 Penulisan Ilmiah}{Kuis}{SCPMK0401-03002}{Kuis untuk mengukur pemahaman struktur karya tulis ilmiah, rumusan masalah, dan metodologi penulisan.}{Menjawab soal pilihan ganda dan/atau uraian singkat tentang penulisan ilmiah.}{Jawaban tertulis kuis.}{Pemahaman struktur karya ilmiah & 3 & Struktur karya ilmiah (abstrak, pendahuluan, tinjauan pustaka, metodologi, pembahasan, kesimpulan) dijelaskan dengan benar dan lengkap. & Sebagian besar struktur dipahami, tetapi beberapa komponen belum dijelaskan dengan tepat. & Pemahaman struktur karya ilmiah tidak memadai atau komponen penting tidak diketahui. \\ \\
Ketepatan rumusan masalah dan metodologi & 2 & Rumusan masalah operasional dan metodologi yang sesuai dapat diidentifikasi dengan benar pada contoh yang diberikan. & Rumusan masalah atau metodologi cukup tepat, tetapi konsistensinya belum terjaga. & Rumusan masalah tidak operasional atau metodologi tidak sesuai dengan topik penelitian. \\}

\assessmentblock{PBL Karya Tulis Ilmiah}{Project Based Learning}{SCPMK0401-03002}{Mahasiswa menyusun karya tulis ilmiah sederhana secara mandiri/kelompok sesuai kaidah akademik.}{Memilih topik, merumuskan masalah, melakukan tinjauan pustaka, menulis, menyunting, dan mempresentasikan.}{Karya tulis ilmiah lengkap dan/atau bahan presentasi.}{Kelengkapan komponen karya tulis ilmiah & 8 & Semua komponen (abstrak, pendahuluan, tinjauan pustaka, metodologi, pembahasan, kesimpulan, daftar pustaka) tersedia dan proporsional. & Sebagian besar komponen tersedia, tetapi beberapa kurang dikembangkan atau tidak proporsional. & Banyak komponen penting tidak ada atau sangat minim pengembangannya. \\ \\
Kualitas isi dan argumentasi ilmiah & 8 & Isi berbasis data/literatur, argumentasi logis, analisis mendalam, dan kesimpulan menjawab rumusan masalah. & Isi cukup berbasis literatur, tetapi argumentasi atau analisis belum sepenuhnya kuat dan terhubung. & Isi tidak berbasis literatur, argumentasi lemah, atau kesimpulan tidak menjawab rumusan masalah. \\ \\
Penulisan ilmiah dan format sitasi & 6 & Bahasa ilmiah baku, sitasi konsisten (APA/IEEE), dan format penulisan sesuai panduan yang ditetapkan. & Bahasa sebagian besar baku, tetapi sitasi atau format ada ketidakkonsistenan yang perlu diperbaiki. & Bahasa tidak ilmiah, sitasi tidak ada atau salah format, atau penulisan tidak sesuai panduan. \\}

\assessmentblock{UAS Presentasi dan Pengumpulan Karya Tulis}{UAS presentasi dan pengumpulan}{SCPMK0301-03001, SCPMK0401-03002}{Mahasiswa mempresentasikan karya tulis ilmiah yang telah disusun dan mengumpulkan versi finalnya.}{Menyiapkan bahan presentasi, mempresentasikan, menjawab pertanyaan, dan mengumpulkan karya tulis final.}{Bahan presentasi, karya tulis final, dan jawaban tanya-jawab.}{Kualitas presentasi lisan & 8 & Presentasi jelas, terstruktur, bahasa lisan efektif, dan mampu menjelaskan isi karya tulis dengan meyakinkan. & Presentasi cukup jelas, tetapi struktur atau kelancaran berbicara masih perlu ditingkatkan. & Presentasi tidak terstruktur, sulit dipahami, atau tidak dapat menjelaskan isi karya tulis. \\ \\
Kualitas karya tulis final & 9 & Karya tulis final lengkap, telah direvisi berdasarkan umpan balik, bahasa baku, dan memenuhi semua kriteria ilmiah. & Karya tulis cukup lengkap, tetapi revisi atau kualitas bahasa masih belum optimal. & Karya tulis tidak lengkap, tidak direvisi, atau banyak kesalahan bahasa yang tidak diperbaiki. \\ \\
Kemampuan menjawab pertanyaan & 6 & Pertanyaan dijawab dengan tepat, menunjukkan pemahaman mendalam terhadap isi karya tulis dan metodologi. & Sebagian besar pertanyaan dijawab dengan cukup baik, tetapi beberapa jawaban kurang lengkap atau tepat. & Tidak mampu menjawab pertanyaan dengan memadai atau menunjukkan pemahaman yang sangat terbatas. \\}

\begin{tblr}{width=\textwidth,colspec={|Q[l,m,3.2cm]|X[l,m]|},hlines,vlines,hspan=minimal,row{1-3}={bg=headgray},rowsep=1.5pt,colsep=3pt}
Jadwal Pelaksanaan: & Minggu 1--17 sesuai RPS. Setiap evaluasi memiliki RTM dan rubrik multi-kriteria. \\
Lain-Lain Yang Diperlukan: & LMS, dokumen penulisan, perangkat pengolah kata, dan perangkat presentasi. \\
Pustaka: & Mengacu pustaka utama dan pendukung pada bagian RPS. \\
\end{tblr}
